\documentclass[10pt]{article}
\usepackage[a4paper,margin=2cm]{geometry}
\usepackage[T1]{fontenc}
\usepackage[utf8]{inputenc}
\usepackage{tgpagella,mathpazo,booktabs,longtable,array}
\renewcommand{\sfdefault}{qpl}
\usepackage[hidelinks]{hyperref}
\newcommand{\pending}[1]{\textbf{[Confirm: #1]}}
\begin{document}
\section*{Supplementary methods: treatment and outcome ascertainment}
\noindent\textit{Implementation-grounded draft.} This description follows the
R scripts supplied in \texttt{opera/R}. It describes observed post-index
events, not causal attribution of events to treatment. External DALY-CARE
codebooks and grading functions must be archived with the final protocol.
Production and sensitivity branches are distinguished below; unresolved
choices are not presented as finalized study methods.

\subsection*{Population, index date and temporal splits}
The supervised prediction population comprises patients with a reconstructed
first-line treatment index and endpoint-specific observation eligibility.
It is distinct from the broader population used for encoder pretraining.
Encoder pretraining used records before 1 January 2022 from the broader patient
population, including patients subsequently assigned to downstream tuning and
held-out sets, according to the study team's account. The supplied pretraining
configuration also specifies this exclusive calendar boundary. This is a
temporal evaluation of subsequent treatment episodes, not a patient-disjoint
evaluation of an encoder that has never encountered these patients.
Outcome files contain one row per eligible patient with patient identifier,
split, index date, first qualifying event date and censoring date. Absence from
an endpoint file denotes masked supervision, not an event-free observation.
The outcome script assigns training before 1 January 2022, tuning during
2022, and held-out from 1 January 2023 using the treatment index date.

\subsection*{First-line episode reconstruction}
Treatment records combine medication/procedure codes, descriptive regimen
text, indication and source information. Records labelled \texttt{SP beh}
are excluded by the supplied code. Hematological indication codes take
priority in the indication filter; non-hematological code/text and drug-text
exclusions screen potentially unrelated treatment. Uncertain records can
remain candidates under the fallback rules.

In the primary mapped-table branch, candidates are searched from 30 days
before the earliest recorded diagnosis to ten years after it. Consecutive
treatment records separated by no more than 21 days form an episode. This is
a consecutive-gap rule, not a maximum total episode duration. Episode evidence
includes anchor drugs, registry source, distinct dates and distinct codes.
Anchor-containing episodes qualify directly. Trusted registry episodes qualify
unless they contain prophylaxis-suspect evidence with fewer than two codes;
repeated anti-CD20 monotherapy can qualify in this branch. Other
prophylaxis-suspect episodes require at least three dates and two codes;
remaining episodes require at least two dates or two codes.

Curated CLL and MM line-one records provide seed dates. Episodes within 30
days of a seed interval are preferentially selected and retained even when
they do not pass the ordinary therapy-like gate. The index remains the
selected episode's start, not the clinical seed date itself. Unseeded patients
receive the earliest qualifying episode. The seed-free comparison is an
ascertainment audit rather than independent validation of the seeded output.

\subsection*{Disease assignment and rescue}
Diagnoses are mapped through the DALY-CARE ICD-10 dictionaries and explicit
supplementary patterns. Registry diagnoses are assigned the highest priority.
Primary assignment uses diagnoses up to 20 years before treatment and a
14-day forward grace. Confirmation accepts qualifying high-priority evidence
(threshold 3 in the call) or repeated evidence separated by at least 14 days.
Fine labels are preferred to broad labels, with specified aggressive-disease
precedence and priority/date/specificity tie-breaking. Derived Richter and
follicular transformations require antecedent CLL/SLL or FL at least 90 days
before the aggressive diagnosis; these are algorithmic labels, not uniformly
pathology-confirmed transformations.

For patients with an index but no assigned cohort, a second pass searches
their diagnosis history for the earliest confirmable specific disease,
re-identifies treatment with a 90-day pre-diagnosis grace and 20-year search,
and assigns disease using a 90-day forward grace and ten-year lookback.
This explicitly permits retrospective coding-lag rescue.
The study team specifies these retrospective identities for cohort analysis,
not prediction-time inputs. Disease-conditioned router models nevertheless
consume disease indices; their label provenance requires separate verification
before extending this statement to those experiments.
Remaining patients may be assigned a cohort from recognizable treatment-arm
prefixes, with provenance \texttt{arm\_rescue}. These labels can have grouped
rather than genuinely fine resolution. \pending{include arm rescue in the
main cohort or sensitivity analysis only.}

\subsection*{Treatment profiles and production branches}
Regimen text and coded treatment evidence are normalized and matched to
disease-specific rules to generate primary/secondary arm labels and exclusion
flags. Episode-wide \texttt{arm\_primary} is the descriptive treatment profile;
day-one arms are separately generated and should not silently replace it.
Unresolved and unmapped labels are retained as ascertainment uncertainty,
not interpreted as confirmed regimens. Cohorts labelled
\texttt{EXCLUDE\_SECONDARY} are removed from the mapped/expanded output.

The supplied script subsequently recomputes \texttt{FROZEN} first-line objects
with anti-CD20 monotherapy disabled, whereas the earlier mapped CSV branch
enables it. Both downstream outcome scripts explicitly reload the mapped CSV.
A later audit constructs a registry-date fallback but does not export it.
\pending{server export provenance agrees with this supplied script revision.}
These branches must not be conflated into one ascertainment algorithm or flow.

\subsection*{Laboratory grading and observation eligibility}
Laboratory values are cleaned with source unit information and analyte
dictionaries. Graded endpoints use the first qualifying measurement strictly
after the index and no later than administrative observation end, usually at
grade $\geq2$ and $\geq3$. Baseline abnormalities are not universally excluded;
this construction is first post-index documentation, not necessarily new onset
or worsening from a normal baseline.

The script contains v6-labelled local graders, v5-labelled fibrinogen logic
whose LLN-fraction structure is also present in v6, and externally supplied
graders obtained from the DALY-CARE package. The inspected upstream revision
\texttt{a0742b0} explicitly references CTCAE v5 for the existing laboratory
graders. For example, its platelet grade-3/4 boundary is
$25\times10^9$/L, whereas CTCAE v6 uses $10\times10^9$/L. The exact
revision loaded on the server remains to be verified.
CTCAE v6.0 was published on 22 July 2025
\footnote{National Cancer Institute, CTCAE v6.0:
\url{https://dctd.cancer.gov/research/ctep-trials/for-sites/adverse-events/ctcae-v6.pdf}.}.
These endpoints represent first post-index documentation of CTCAE-based
laboratory abnormalities, not adjudicated treatment-attributable toxicity.
The final endpoint dictionary must identify the reference version, numerical
thresholds, units, reference limits, baseline handling and approximations for
each endpoint rather than claiming universal v6 compliance.
For example, local hyponatremia grading omits symptom-based escalation;
hyperuricemia uses a laboratory-only elevation rather than adjudicated
physiological consequences; LDH is a binary elevation endpoint.
Fibrinogen grading omits the baseline-relative axis and uses unit-sensitive
fallback and absolute-floor constants requiring source-unit verification.
Local interval boundaries also require validation against reference inequalities.
The file named \texttt{leukocytosis\_g2plus} requests grade 3; CTCAE v6 has no
grade-2 leukocytosis category and advises against using the term for increases
due to leukemia, steroids or growth factors. This endpoint is therefore
exploratory disease-state prediction, not a primary toxicity claim.
The local hypokalemia bins are not CTCAE v6: v6 grade 2 depends on symptoms,
and its laboratory grade-3 threshold differs from the local historical bins.
The externally derived graded lymphocyte-decrease endpoints reflect a v5 term;
v6 lists lymphopenia only as present/grade 1. These labels require a versioned
crosswalk before a uniform-v6 analysis is claimed.

The active coverage engine uses any recorded laboratory activity, not a
separate coverage window for each analyte. Default requirements are a draw in
the 365 days up to and including index and a draw from index through day 30.
Source coverage is index on/after 1 January 2017 or recorded laboratory
activity in the observation window. Early observed death can exempt the entry
draw requirement after go-live; the optional pre-go-live exemption is disabled
by default. Observing an event does \emph{not} override these eligibility rules,
despite an older comment saying otherwise.

For eligible patients without an event, censoring is the earlier of the last
post-index draw plus 180 days and administrative observation end. For eligible
patients with an event, censoring is administrative observation end. Because
coverage is shared across analytes, this documents monitoring rather than
proving that each endpoint-specific test was performed throughout follow-up.

\subsection*{Other endpoints and source-specific coverage}
Procedure endpoints use the first qualifying post-index coded event.
Diagnosis-based incident endpoints exclude qualifying ICD prefixes observed
on/before index (or the endpoint's configured grace), whereas acute infection
diagnoses use their dedicated first-post-index construction. Medication-defined
infection is an antimicrobial-use proxy, not microbiologically confirmed
infection: the implemented IV branch uses SMR administrations, with a minimum
one-day episode span and seven-day separation. Prescription-based events
are built separately and require dispensing dates spanning at least one day;
this is not measured days of supply, and a single dispensing date cannot qualify.
Their calls currently use the shared coverage-gated
wrapper, so laboratory eligibility also applies. Prophylaxis filtering is not
identical between these branches and requires explicit review.

The source gate excludes designated LPR3-dependent endpoints indexed before
3 February 2019 plus a 30-day buffer (5 March 2019), and IV-infection endpoints
before the configured SMR date of 15 May 2018. These are implemented analysis
eligibility dates, not a claim that national capture became complete on those
dates. Ungated backups and audit files are not primary endpoints.
\pending{confirm these coverage dates against source diagnostics.}

\subsection*{Second-line treatment and composites}
Episode-derived second-line ascertainment is disease-group-specific. It uses
first-line end, minimum post-end gaps, consolidation overlap, maintenance-only
codes and transplantation guards; single-day first-line episodes can have
their end imputed. Qualifying second-line events require at least moderate
evidence, with registry support, anchor/indication and corroborating treatment
evidence contributing to confidence. Weak-evidence-only patients are masked,
not made negative. Low-reliability BL/LBL and HCL episode-derived observations
are masked, with registry-confirmed rescue handled separately. CLL/MM use
curated line-two records rather than the same episode algorithm.
This is time to next qualifying treatment, not directly observed progression.

Composite eligibility is the intersection of component populations. Its event
is the earliest component event within the earliest component censoring date.
Treatment failure combines second-line treatment and death; the expanded
failure endpoint adds the next hematological malignancy definition.
Severe cytopenia combines grade-3 neutropenia, thrombocytopenia and anemia;
hepatotoxicity combines grade-3 ALT, AST and bilirubin; the all-definitions
infection composite combines serious infection, prescription and IV proxies.
Separate co-occurrence endpoints require abnormalities on multiple axes
within specified windows and must not be described as these simple unions.
The TLS proxy requires two of four axes within seven days, confined to the first
seven post-index days. Hy's-law-style liver injury requires transaminase and
bilirubin axes within 14 days, without full cholestasis or alternative-cause
exclusion. DIC-like coagulopathy requires two of three axes within seven days;
it is not ISTH-scored DIC. The supplied co-occurrence constructor does not apply
the ordinary laboratory monitoring mask or last-draw censoring rule, and these
three endpoints are absent from the source-gate list.
\pending{adequate observation requirements for co-occurrence negatives.}
Death is a primary event when overall survival is a component; otherwise
competing death is recovered downstream using patient death information.

\subsection*{Follow-up, versioning and flow counts}
The main outcome script specifies an administrative cutoff of 19 March 2026;
the second-line script specifies 4 March 2026. \pending{reconcile or explicitly
document endpoint-specific follow-up dates.} Status distinguishes actual death
from the combined death/follow-up date. Patients lacking valid follow-up after
their treatment index must be accounted for separately from treatment
ascertainment failure.

First-line instrumentation writes exact internal aggregate counts for records,
episodes and unique patients at treatment filtering, episode qualification,
first-pass selection, diagnosis rescue, arm rescue and CSV export. Branches and
counting units are explicit; marginal breakdowns are not mutually exclusive
exclusions. Downstream no-follow-up exclusions, endpoint eligibility and temporal
split counts must be joined from their own pipeline stages before populating
the final flow diagram. Small-cell and complementary-disclosure review applies
to the full publication bundle.
The coverage gate now leaves existing exclusion-audit files intact when a
repeated run finds no pre-coverage patients. Original audit artifacts should
still be archived with the export provenance for flow-count verification.
\end{document}
